% \chapter{ESP-IDF Build Instructions}

% \section{Prerequisites}

% \subsection{Software Requirements}
% \begin{itemize}
%     \item ESP-IDF v5.0 or later
%     \item Python 3.8 or later
%     \item CMake 3.16 or later
%     \item Ninja build system
%     \item Git
% \end{itemize}

% \subsection{Hardware Requirements}
% \begin{itemize}
%     \item ESP32-S3-WROOM-1 Development Kit
%     \item OV3660 camera module
%     \item USB cable for programming
%     \item 3.7V LiPo battery (optional, for power testing)
% \end{itemize}

% \section{ESP-IDF Installation}

% \subsection{Linux / macOS}

% \begin{lstlisting}[language=bash]
% # Clone ESP-IDF repository
% git clone --recursive https://github.com/espressif/esp-idf.git
% cd esp-idf
% git checkout v5.0

% # Install dependencies
% ./install.sh esp32s3

% # Set up environment variables
% . ./export.sh
% \end{lstlisting}

% \subsection{Windows}

% Download and run the ESP-IDF installer from:\\
% \texttt{https://dl.espressif.com/dl/esp-idf/}

% \section{Project Setup}

% \begin{lstlisting}[language=bash]
% # Clone the project repository
% git clone [your-repository-url]
% cd edge-ai-vineyard-monitoring/esp-detection/esp-detection

% # Configure the project
% idf.py menuconfig
% \end{lstlisting}

% \subsection{Key Configuration Options}

% Navigate through menuconfig:
% \begin{itemize}
%     \item \textbf{Component config → ESP32S3-Specific}
%         \begin{itemize}
%             \item Support for external SPIRAM: Enable
%             \item SPI RAM config → Mode: Octal
%             \item SPI RAM config → Speed: 80MHz
%         \end{itemize}
%     \item \textbf{Serial flasher config}
%         \begin{itemize}
%             \item Flash size: 16MB
%             \item Flash SPI mode: QIO
%         \end{itemize}
%     \item \textbf{Compiler options}
%         \begin{itemize}
%             \item Optimization Level: -O2
%         \end{itemize}
% \end{itemize}

% \section{Building the Project}

% \begin{lstlisting}[language=bash]
% # Clean previous builds (optional)
% idf.py fullclean

% # Build the project
% idf.py build
% \end{lstlisting}

% The build process generates:
% \begin{itemize}
%     \item \texttt{bootloader.bin} - First-stage bootloader
%     \item \texttt{partition-table.bin} - Partition table
%     \item \texttt{app-name.bin} - Application binary
% \end{itemize}

% \section{Flashing the Device}

% \subsection{Automatic Flashing}

% \begin{lstlisting}[language=bash]
% # Flash and monitor in one command
% idf.py -p /dev/ttyUSB0 flash monitor
% \end{lstlisting}

% Replace \texttt{/dev/ttyUSB0} with your serial port:
% \begin{itemize}
%     \item Linux: \texttt{/dev/ttyUSB0} or \texttt{/dev/ttyACM0}
%     \item macOS: \texttt{/dev/cu.usbserial-*}
%     \item Windows: \texttt{COM3} (or check Device Manager)
% \end{itemize}

% \subsection{Manual Flashing with esptool}

% \begin{lstlisting}[language=bash]
% esptool.py --chip esp32s3 --port /dev/ttyUSB0 --baud 921600 \
%   write_flash -z \
%   0x0 build/bootloader/bootloader.bin \
%   0x8000 build/partition_table/partition-table.bin \
%   0x10000 build/app-name.bin

% # Verify flash
% esptool.py --chip esp32s3 --port /dev/ttyUSB0 verify_flash \
%   0x10000 build/app-name.bin
% \end{lstlisting}

% \section{Serial Monitoring}

% \begin{lstlisting}[language=bash]
% # Monitor serial output
% idf.py -p /dev/ttyUSB0 monitor

% # Exit monitor: Ctrl + ]
% \end{lstlisting}

% \subsection{Expected Output}

% \begin{lstlisting}[language=bash]
% I (328) main: ESP32-S3 Vineyard Monitor starting...
% I (338) camera: Initializing OV3660 camera
% I (445) camera: Camera initialized successfully
% I (446) psram: PSRAM initialized: 8192 KB
% I (450) main: Loading detection model...
% I (1250) detection: ESPDet-Pico loaded (478 KB)
% I (1255) main: Loading classification model...
% I (3890) classification: MobileNetV2 loaded (2.73 MB)
% I (3895) main: System ready. Starting inference...
% I (3978) camera: Frame captured (10245 bytes)
% I (4344) detection: Detected 8 leaves in 366ms
% I (10189) classification: Classified 8 leaves in 5845ms
% I (10195) aggregation: Final prediction: Healthy (98.7%)
% \end{lstlisting}

% \section{Common Issues and Solutions}

% \begin{table}[H]
% \centering
% \caption{Troubleshooting common build/flash issues}
% \begin{tabular}{@{}p{4cm}p{8cm}@{}}
% \toprule
% \textbf{Issue} & \textbf{Solution} \\
% \midrule
% Failed to connect to device & Press and hold BOOT button, then press RESET. Release RESET, then release BOOT \\
% \midrule
% Out of memory during build & Enable PSRAM in menuconfig, increase optimization level \\
% \midrule
% Camera initialization failed & Check camera connections, verify camera I2C address (0x3C), ensure sufficient power supply \\
% \midrule
% Model file not found & Ensure .espdl model files are in \texttt{models/} directory, check CMakeLists.txt \\
% \midrule
% Brownout detector triggered & Use external 5V power supply instead of USB, check power cable quality \\
% \bottomrule
% \end{tabular}
% \end{table}

% \section{OTA (Over-The-Air) Updates}

% \subsection{Partition Table Configuration}

% Edit \texttt{partitions.csv}:
% \begin{lstlisting}
% # Name,   Type, SubType, Offset,  Size
% nvs,      data, nvs,     0x9000,  0x6000
% phy_init, data, phy,     0xf000,  0x1000
% factory,  app,  factory, 0x10000, 2M
% ota_0,    app,  ota_0,   ,        2M
% ota_1,    app,  ota_1,   ,        2M
% ota_data, data, ota,     ,        0x2000
% \end{lstlisting}

% \subsection{OTA Update Code}

% \begin{lstlisting}[language=C++]
% #include "esp_https_ota.h"

% esp_http_client_config_t config = {
%     .url = "https://example.com/firmware.bin",
%     .cert_pem = server_cert_pem_start,
% };

% esp_err_t ret = esp_https_ota(&config);
% if (ret == ESP_OK) {
%     esp_restart();
% }
% \end{lstlisting}

% \section{Performance Profiling}

% \subsection{Memory Profiling}

% \begin{lstlisting}[language=C++]
% #include "esp_heap_caps.h"

% // Check free PSRAM
% size_t free_psram = heap_caps_get_free_size(MALLOC_CAP_SPIRAM);
% ESP_LOGI(TAG, "Free PSRAM: %d bytes", free_psram);

% // Check largest free block
% size_t largest = heap_caps_get_largest_free_block(MALLOC_CAP_SPIRAM);
% ESP_LOGI(TAG, "Largest free block: %d bytes", largest);
% \end{lstlisting}

% \subsection{CPU Timing}

% \begin{lstlisting}[language=C++]
% #include "esp_timer.h"

% int64_t start = esp_timer_get_time();
% // Your code here
% int64_t end = esp_timer_get_time();
% ESP_LOGI(TAG, "Execution time: %lld us", end - start);
% \end{lstlisting}

% \section{Additional Commands}

% \begin{lstlisting}[language=bash]
% # Erase entire flash
% idf.py -p /dev/ttyUSB0 erase-flash

% # Read partition table
% esptool.py --chip esp32s3 --port /dev/ttyUSB0 read_flash \
%   0x8000 0x1000 partition_table.bin

% # Monitor with filtering
% idf.py monitor | grep -E "detection|classification|aggregation"

% # Core dump analysis
% idf.py -p /dev/ttyUSB0 coredump-info

% # GDB debugging
% idf.py -p /dev/ttyUSB0 gdb
% \end{lstlisting}

% \section{Useful Resources}

% \begin{itemize}
%     \item ESP-IDF Programming Guide: \url{https://docs.espressif.com/projects/esp-idf/}
%     \item ESP32-S3 Technical Reference Manual: \url{https://www.espressif.com/sites/default/files/documentation/esp32-s3_technical_reference_manual_en.pdf}
%     \item ESP-DL Documentation: \url{https://github.com/espressif/esp-dl}
%     \item Espressif Forum: \url{https://esp32.com/}
% \end{itemize}
